This section documents the search for evidence that could test the decision of
Section~\ref{sec:decision}, so that the absence of a shipboard evaluation is a
recorded finding rather than an assertion. The outcome is negative: no source
examined can test the model, and the one source retained serves only as physical
context.

To estimate~\eqref{eq:model} and to test a plan, a source must provide, for each
population, an installed count, an exposure over a bounded horizon, a count of
failures in that horizon distinguished from preventive replacements, and a
data-holder label. A source missing any one of the four cannot test the decision,
and no amount of modelling supplies it. Of the \nSources{} sources examined,
\nSuitable{} meet that requirement, and the requirement each fails is recorded in
the accompanying audit file.



No public source was found giving failures of shipboard power-electronic equipment
with installed populations and exposure. Port-state-control databases record
inspection deficiencies per ship, which are findings against a standard rather
than failures of an identified population, and marine casualty databases record
incidents without a denominator; neither supports a rate. The closest public
analogue is a converter fleet with a site structure, and the largest openly
licensed candidate was examined: operations, maintenance, weather and production
data for 174 anonymised utility-scale photovoltaic plants~\cite{oedi_pv}. It was
rejected on inspection, since its maintenance fields record weather-event tickets
and an affected-asset type rather than inverter replacement counts, with no
inverter population per site and no exposure. The accelerated-ageing sets for
power devices were also examined; the IGBT set comprises six devices characterised
at waveform level, with no populations, no in-service exposure and no data
holders, and was rejected.

\subsection{Physical Context: Capacitor Degradation}
Aluminium electrolytic capacitors are the usual life-limiting element of the
converter and power-supply stages considered here, and their wear-out under
electrical and thermal stress is documented~\cite{celaya2012,kulkarni2012}.
The retained set comprises \capUnits{} such capacitors held at
\capLevels\,V bias for \capHours\,h, with impedance sweeps giving equivalent
series resistance (ESR) and capacitance throughout.

Final ESR relative to each unit's own initial value rises monotonically with
stress, with means of \capFinalTen, \capFinalTwelve{} and \capFinalFourteen{} at
10, 12 and 14\,V and a within-level standard deviation of \capSd, so a stress
effect and unit-to-unit variability are both visible in real power-electronic
components.

One fact limits what may be concluded, and is reported because the tabulation
distributed with such data can mislead: the source's failure criterion is an ESR
of \capCrit{} times initial; applied sweep by sweep, \capAnyCross{} of the
\capUnits{} units register a crossing, but every crossing occurs in the first few
sweeps and is followed by a decline, and \capEndsAbove{} units end the test above
the criterion, the largest final ratio being \capMaxFinal. The crossings are
transients of the early measurements rather than threshold crossings sustained to
end of life, so this set yields no credible times to failure and no failure rate
is taken from it.

The set therefore supports two statements only, neither validating the spares
model: the assumed degradation mechanism is real and measurable, and even under
accelerated stress such a population can run for thousands of hours with no
confirmed failure. It contains no shipboard population, no mission, no estimated
rate and no plan compared with observed failures.


